%% LyX 2.5.1 created this file.  For more info, see https://www.lyx.org/.
%% Do not edit unless you really know what you are doing.
\documentclass[english,hebrew]{article}
\usepackage[HE8,T1]{fontenc}
\usepackage[utf8]{inputenc}
\usepackage{babel}
\usepackage{amssymb}
\usepackage[unicode=true]
 {hyperref}
\begin{document}
\title{מבוא לטופולוגיה קבוצתית: מרחבים טופולוגיים, בסיסים, וטופולגיית הסדר}
\maketitle
\begin{description}
\item [{קטגוריות:}] טופולוגיה
\item [{תגים:}] מרחבים מטריים, מרחבים טופולוגיים
\item [{מזהה:}] \L{topological\_spaces\_and\_bases}
\end{description}

\section{בסיסים ותת-בסיסים}

בפוסטים הקודמים דיברתי על תכונות של קבוצות פתוחות וסגורות במרחבים
מטריים, במטרה לקבל תחושה שמתקיימים שני דברים מרכזיים:
\begin{itemize}
\item הרבה ממה שאנחנו עושים במרחבים מטריים אפשר לנסח בלשון של קבוצות פתוחות
בלי להזדקק לדיבורים על המטריקה.
\item יש סט תכונות פשוט למדי שמאפיין קבוצות פתוחות בלי להזדקק לדיבורים על
המטריקה.
\end{itemize}
כלומר, יש לנו הזדמנות לבצע \textbf{אבסטרטקציה} של מה שהולך במרחבים
מטריים ולתפוס מחלקה רחבה בהרבה של מרחבים, שעדיין יהיו מעניינים בצורה
דומה לזו שבה מרחבים מטריים מעניינים - \textbf{מרחבים טופולוגיים}.
ההגדרה שאני אציג כאן, שמסתמכת על התכונות שראינו בפוסטים הקודמים, היא
לא איזו אמת קוסמית צרופה; זו פשוט ההגדרה שהשתרשה בעולם המתמטי אחרי
שבגלגולים המוקדמים של הטופולוגיה הקבוצתית בתחילת המאה ה-{\beginL 20\endL}
כל אחד מהחלוצים עבד עם סט הגדרות משל עצמו.

\textbf{מרחב טופולוגי} הוא זוג \L{$\left(X,\mathcal{T}\right)$} של
קבוצה \L{$X$} שהיא ה\textquotedblright עולם\textquotedblright{} שבו
נמצאים האיברים שלנו, שעליהם אנחנו חושבים בתור \textquotedblleft נקודות\textquotedblright{}
)מכאן - \L{Point set topology}, השם באנגלית של \textquotedblleft טופולוגיה
קבוצתית\textquotedblright ( ו-\L{$\mathcal{T}\subseteq\mathcal{P}\left(X\right)$}
הוא אוסף של תת-קבוצות של \L{$X$} שאנחנו קוראים להן \textbf{קבוצות
פתוחות}. הדרישה מ-\L{$\mathcal{T}$} היא שיתקיימו שלוש התכונות הבאות:
\begin{itemize}
\item \L{$\emptyset,X\in\mathcal{T}$}
\item אם יש לנו אוסף קבוצות \L{$A_{i}\in\mathcal{T}$} אז \L{$\bigcup_{i}A_{i}\in\mathcal{T}$}
\item אם יש לנו מספר סופי של קבוצות \L{$A_{1},\ldots,A_{n}\in\mathcal{T}$}
אז \L{$\bigcap^{n}_{i=1}A_{i}\in\mathcal{T}$}
\end{itemize}
התכונות הללו הן \textbf{האקסיומות} של מרחב טופולוגי. שימו לב שהכוונה
לאקסיומות במובן המודרני של המילה, כלומר \textbf{לא} משהו כמו \textquotedblleft טענות
שהנכונות שלהן מובנת מאליה ולא צריך להוכיח אותה\textquotedblright{}
אלא \textquotedblleft סט ספציפי של \textbf{הנחות} שאנחנו מצפים מאובייקט
לקיים כדי שיהיה מעניין\textquotedblright , בדומה לאקסיומות של \textbf{חבורה}
או של \textbf{מרחבים וקטוריים}.

נתנו הגדרה? יופי, הדבר הראשון שצריך לתת אחרי הגדרה הוא \textbf{דוגמאות}
למרחבים שעונים להגדרה הזו. אבל כאן אנחנו טיפה מסתבכים, כי לתת את \L{$\mathcal{T}$}
במפורש יכול להיות מעייף מאוד. בואו נחשוב לדוגמא על מה שקרה במרחב הטופולוגי
הפשוט ביותר שראינו - המרחב המטרי \L{$\mathbb{R}$}. שם היו לנו קבוצות
פתוחות שהן \textbf{כדורים פתוחים}, מה שהיה בסך הכל שם אחר לקטעים פתוחים
מהצורה \L{$\left(a,b\right)$} )רק שהצגתי אותם בתור \L{$\left(a-r,a+r\right)$}
עבור \L{$a\in\mathbb{R}$} ו-\L{$r>0$} ממשיים(. אבל גם איחודים של
הקטעים הפתוחים הללו הם קבוצות פתוחות, אפילו איחודים אינסופיים, וגם
אפשר לחתוך קבוצות... אין לי תיאור פשוט של הקבוצות הפתוחות אפילו במרחב
הזה. אבל שימו לב מה כן יש לי - את הדבר הזה שהתחלתי ממנו, כדורים פתוחים.
זה מוביל אותנו לרעיון של \textbf{בסיס} לטופולוגיה: בסיס הוא אוסף של
קבוצות שבעזרתן אפשר לתאר את כל הטופולוגיה, ולכן כשרוצים להגדיר טופולוגיה
מספיק לתאר בסיס שלה ולא את כל הקבוצות הפתוחות במפורש. זה מקל בצורה
משמעותית על הגדרת טופולוגיות.

איך בדיוק נגדיר בסיס? גם פה המרחבים המטריים באים לעזרתנו. במרחבים
הללו הקבוצות הפתוחות ה\textquotedblright בסיסיות\textquotedblright{}
היו הכדורים הפתוחים; למעשה, \textbf{הגדרנו }מה זו קבוצה פתוחה בעזרתם.
אמרנו ש-\L{$A$} היא קבוצה פתוחה אם לכל \L{$a\in A$} קיים כדור פתוח
\L{$B\left(a,r\right)\subseteq A$} עם \L{$r>0$}. מסקנה אחת שאפשר
להסיק מייד ואולי לא כזו ברורה ממבט ראשון היא זו: אם \L{$A$} היא קבוצה
פתוחה \textbf{כלשהי} וניקח אוסף של כדורים פתוחים \L{$B\left(a,r_{a}\right)\subseteq A$}
לכל \L{$a\in A$}, ואז נסתכל על האיחוד שלו, נקבל ש-

\L{$\bigcup_{a\in A}B\left(a,r_{a}\right)\subseteq A$}

כי איחוד של קבוצות שכולן מוכלות ב-\L{$A$} מוכל ב-\L{$A$}. מצד שני,
לכל \L{$a\in A$} מתקיים \L{$a\in B\left(a,r_{a}\right)$} ולכן

\L{$A\subseteq\bigcup_{a\in A}B\left(a,r_{a}\right)$}

ומשני אלו נסיק \L{$A=\bigcup_{a\in A}B\left(a,r_{a}\right)$}. במילים
אחרות - במרחב מטרי אפשר לתאר \textbf{כל} קבוצה פתוחה בתור איחוד )אולי
אינסופי( של כדורים פתוחים. 

אם כן, התכונה הזו של כדורים פתוחים - שלכל קבוצה פתוחה \L{$A$} וכל
\L{$a\in A$} קיים כדור פתוח שמכיל את \L{$a$} ומוכל ב-\L{$A$} -
זו תכונה טובה; זה מה שהיינו רוצים לקבל גם באופן כללי. כלומר - היינו
רוצים שבסיס יהיה קבוצה \L{$\mathcal{B}$} כלשהי כך שלכל \L{$A\in\mathcal{T}$}
וכל \L{$a\in A$} יהיה קיים \L{$B\in\mathcal{B}$} כך ש-\L{$a\in B\subseteq A$}.
בפרט, היינו רוצים שאפשר יהיה \textbf{להגדיר} את \L{$\mathcal{T}$}
ככה: בהינתן \L{$\mathcal{B}$}, נגדיר את הטופולוגיה \L{$\mathcal{T}$}
שנוצרת בעזרתו בתור אוסף כל הקבוצות \L{$A\subseteq X$} שעבורן לכל
\L{$a\in A$} קיים \L{$B\in\mathcal{B}$} כך ש-\L{$a\in B\subseteq A$}.

אבל - האם אפשר סתם לקחת \textbf{כל} אוסף \L{$\mathcal{B}$} של קבוצות
ולהגדיר ככה \L{$\mathcal{T}$} בעזרתו? כמובן שלא, אנחנו רוצים שהאקסיומות
של מרחב טופולוגי יתקיימו. בואו ננסה להוכיח אותן ונראה מה בעצם צריך
לדרוש מ-\L{$\mathcal{B}$} כדי שזה יעבוד.

בתור התחלה, בשביל \L{$\emptyset\in\mathcal{T}$} לא צריך כלום - התנאי
שנדרש מ-\L{$\emptyset$} כדי להיות פתוחה היא שלכל \L{$a\in\emptyset$}
יהיה איבר בסיס שכך-וכך, אבל אין \L{$a$} שיקיים \L{$a\in\emptyset$}
אז התנאי הזה מתקיים \textbf{באופן ריק}.

שנית, בשביל \L{$X\in\mathcal{T}$} צריך ש\textbf{כל} \L{$x\in X$}
יהיה שייך לקבוצה כלשהי ב-\L{$\mathcal{B}$}. זה בדיוק מה שהתנאי על
\textquotedblleft\L{$X$} פתוחה\textquotedblright{} אומר במקרה הזה:
שלכל \L{$x\in X$} יהיה קיים \L{$B\in\mathcal{B}$} כך ש-\L{$x\in B$}
וגם \L{$B\subseteq X$}, והתנאי השני הזה מתקיים תמיד.

התנאי השני על \L{$\mathcal{T}$}, זה על איחודים, מתקבל בקלות בחינם:
אם \L{$\left\{ A_{i}\right\} _{i\in\Lambda}$} היא משפחה כלשהי של
קבוצות פתוחות ו-\L{$A=\bigcup_{i\in\Lambda}A_{i}$} אז בואו ניקח \L{$a\in A$}
שרירותי כלשהו. מהגדרת איחוד, קיים \L{$i$} כך ש-\L{$a\in A_{i}$}.
מכך ש-\L{$A_{i}$} פתוחה נקבל שקיים \L{$B\in\mathcal{B}$} כך ש-\L{$a\in B\subseteq A_{i}\subseteq A$}
וזה מה שרצינו.

לגבי תכונת החיתוך שווה לזכור שכל מה שאנחנו צריכים בפועל הוא לדבר על
חיתוך של \textbf{שתי} קבוצות, כי אם נוכיח שעבור \L{$A_{1},A_{2}\in\mathcal{T}$}
מתקיים \L{$A_{1}\cap A_{2}\in\mathcal{T}$}, אפשר להוכיח את זה באינדוקציה
לכל מספר סופי של קבוצות: נסתכל על החיתוך הכללי \L{$A_{1}\cap A_{2}\cap\ldots\cap A_{n}$}
ונציג אותו בתור \L{$A_{1}\cap\left(A_{2}\cap\ldots\cap A_{n}\right)$}.
נשתמש בהנחת האינדוקציה כדי לקבל ש-\L{$A_{2}\cap\ldots\cap A_{n}$}
היא קבוצה פתוחה )כי היא חיתוך של \L{$n-1$} קבוצות( ועכשיו נשתמש בכך
שחיתוך של שתי קבוצות פתוחות הוא פתוח.

אם כן, יהיו \L{$A_{1},A_{2}$} שתי קבוצות פתוחות ונסתכל על \L{$A_{1}\cap A_{2}$}.
ניקח \L{$a\in A_{1}\cap A_{2}$}. אנחנו \textbf{רוצים} שיהיה קיים
\L{$B\in\mathcal{B}$} כך ש-\L{$a\in B\subseteq A_{1}\cap A_{2}$}.
עכשיו, מכיוון ש-\L{$a\in A_{1}\cap A_{2}$} אני מקבל שמתקיים גם \L{$a\in A_{1}$}
וגם \L{$a\in A_{2}$}. מכך שאלו קבוצות פתוחות נקבל \L{$B_{1},B_{2}\in\mathcal{B}$}
כך ש-\L{$a\in B_{1}\subseteq A_{1}$} וגם \L{$a\in B_{2}\subseteq A_{2}$},
ועכשיו כדי לקבל את \L{$B$} אנחנו... אה... אוקיי, צריך להכניס כאן
דרישה חדשה מ-\L{$\mathcal{B}$}. מה שאני דורש הוא שלכל שני \textbf{אברי
בסיס} \L{$B_{1},B_{2}$} ולכל \L{$a\in B_{1}\cap B_{2}$} יהיה קיים
איבר בסיס \L{$B$} כך ש-\L{$a\in B\subseteq B_{1}\cap B_{2}$}. ה-\L{$B$}
הזה הוא מה שהיינו צריכים, כי הוא יקיים \L{$B\subseteq B_{1}\cap B_{2}\subseteq A_{1}\cap A_{2}$},
והעובדה שהדרישה מנוסחת מראש על קבוצות ב-\L{$\mathcal{B}$} תהפוך אותה
לקלה יחסית להוכחה )בהינתן בסיס \L{$\mathcal{B}$} פשוט יחסית(.

אז אם לסכם, \textbf{בסיס} \L{$\mathcal{B}$} הוא אוסף של קבוצות שמקיים
\begin{itemize}
\item לכל \L{$x\in X$} קיים \L{$B\in\mathcal{B}$} כך ש-\L{$x\in B$}.
\item אם \L{$B_{1},B_{2}\in\mathcal{B}$} ו-\L{$x\in B_{1}\cap B_{2}$}
אז קיים \L{$B\in\mathcal{B}$} כך ש-\L{$x\in B\subseteq B_{1}\cap B_{2}$}.
\end{itemize}
והטופולוגיה ש-\L{$\mathcal{B}$} יוצר היא בעלת התכונה שלכל קבוצה פתוחה
\L{$A$} ולכל \L{$a\in A$}, קיים איבר בסיס \L{$B_{a}\in\mathcal{B}$}
כך ש-\L{$a\in B_{a}\subseteq A$}. בפרט שימו לב שזה אומר ש-\L{$\bigcup_{a\in A}B_{a}=A$},
כלומר כל קבוצה \L{$A$} היא איחוד של אברי בסיס; אבל זה לא עובד נחמד
כמו עם בסיסים באלגברה לינארית - בהחלט ייתכן שלאותה קבוצה פתוחה יהיו
כמה הצגות שונות בתור איחוד של אברי בסיס.

אם כן, אפשר לחשוב על \textbf{בסיס} בתור \textquotedblleft משפחת קבוצות
\L{$\mathcal{B}$} שכל קבוצה פתוחה ב-\L{$\mathcal{T}$} היא איחוד
)אולי אינסופי( שלהן\textquotedblright . ומה אם היינו קצת יותר מגמישים
את הדרישה ומדברים על \textquotedblleft משפחת קבוצות \L{$\mathcal{S}$}
שכל קבוצה פתוחה ב-\L{$\mathcal{T}$} היא איחוד )אולי אינסופי( של \textbf{חיתוכים}
)סופיים( שלהן\textquotedblright ? הדבר הזה נקרא \textbf{תת-בסיס};
הרעיון הוא שתת-בסיס תמיד מייצר בסיס פשוט על ידי כך שלוקחים חיתוכים
סופיים של האיברים שלו. בואו נראה את זה בפעולה.

ראשית, מה בעצם אוסף \L{$\mathcal{S}$} צריך לקיים כדי להיות תת-בסיס?
נראה שהדרישה הראשונה מבסיסים חייבת להתקיים כאן:
\begin{itemize}
\item לכל \L{$x\in X$} קיים \L{$S\in\mathcal{S}$} כך ש-\L{$x\in S$}.
\end{itemize}
כי הרי אם יש \L{$x$} שלא שייך לאף קבוצה בתת-הבסיס, אבוד לנו; אין
שום דרך שאיחודים של חיתוכים של קבוצות התת-בסיס יתנו לנו את \L{$X$}.
אבל האם צריך לדרוש עוד משהו? בואו נראה. אמרתי שבהינתן תת-בסיס \L{$\mathcal{S}$}
אני יכול להגדיר בסיס \L{$\mathcal{B}$} פשוט על ידי זה שאקח את כל
החיתוכים מהצורה \L{$\bigcap^{n}_{i=1}S_{i}$} עבור כל סדרה \L{$S_{1},\ldots,S_{n}$}
מאורך סופי של אברי \L{$\mathcal{S}$}. מה שצריך הוא להוכיח ש-\L{$\mathcal{B}$}
הוא בסיס ומכיוון ש-\L{$\mathcal{S}\subseteq\mathcal{B}$} ברור שהתנאי
הראשון של בסיס מתקיים עבור \L{$\mathcal{B}$} כי הוא התקיים עבור \L{$\mathcal{S}$}.
ומה עם התנאי השני?
\begin{itemize}
\item אם \L{$B_{1},B_{2}\in\mathcal{B}$} ו-\L{$x\in B_{1}\cap B_{2}$}
אז קיים \L{$B\in\mathcal{B}$} כך ש-\L{$x\in B\subseteq B_{1}\cap B_{2}$}.
\end{itemize}
טוב, זה די ברור שאנחנו מקבלים את זה בחינם, כן? הרי אפשר לכתוב

\L{$B_{1}=S^{1}_{1}\cap\ldots\cap S^{1}_{n}$}

\L{$B_{2}=S^{2}_{1}\cap\ldots\cap S^{2}_{m}$}

ולקבל

\L{$B_{1}\cap B_{2}=\left(S^{1}_{1}\cap\ldots\cap S^{1}_{n}\right)\cap\left(S^{2}_{1}\cap\ldots\cap S^{2}_{m}\right)$}

וגם זה איבר של \L{$\mathcal{B}$} כי גם זה חיתוך סופי של אברי \L{$\mathcal{S}$}
אז להגדיר תת-בסיס זה ממש פשוט - רק צריך לוודא שכל איבר במרחב שייך
לקבוצה כלשהי.

עכשיו אפשר סוף כל סוף לראות דוגמאות לטופולוגיות. כבר נמאס לנו לחלוטין
ממרחבים מטריים )או שאני היחיד?( אבל בכל זאת הם אלו שנתנו לנו את ההשראה
אז בואו נציין במפורש שבמרחבים מטריים הבסיס לטופולגיה הוא \textbf{הכדורים
הפתוחים}. עכשיו בואו נראה משהו חדש.

\section{טופולוגיית הסדר}

כבר אמרתי בפוסט הזה שבסיס לטופולוגיה הרגילה על \L{$\mathbb{R}$} הוא
אוסף הקטעים הפתוחים \L{$\left(a,b\right)$}. את האוסף הזה אפשר היה
לראות בתור משהו שמגיע ממטריקה וכל קטע פתוח כזה הוא בעצם כדור פתוח
וכדומה, אבל אפשר היה גם לוותר על מושג המטריקה לחלוטין מלכתחילה. קטע
פתוח זה לא משהו שדורש מטריקה כדי להגדיר אותו. ההגדרה היא

\L{$\left(a,b\right)=\left\{ x\in\mathbb{R}\ |\ a<x<b\right\} $}

בשביל זה אנחנו צריכים דבר אחד: \textbf{יחס סדר} על \L{$\mathbb{R}$}.
יחס הסדר הזה מקיים תכונה פשוטה אחת - הוא \textbf{לינארי}, שזה שם אחר
ל\textbf{יחס סדר מלא} שזו דרך להגיד שלכל \L{$x,y$} מתקיים \L{$x\le y$}
או \L{$y\le x$}, \textquotedblleft כל שני איברים ניתנים להשוואה\textquotedblright .
יש עוד תכונות מועילות שהוא מקיים, למשל \textbf{אקסיומת החסם העליון},
אבל אין לנו צורך בהן כאן.

המספרים הממשיים הם רק קבוצה אחת. אם יש לנו קבוצה סדורה בסדר מלא כלשהו,
\L{$\left(P,\le\right)$}, האם עדיין ניתן להגדיר עליה טופולוגיה באותה
דרך? לקחת את הבסיס שהוא כל הקבוצות \L{$\left(a,b\right)=\left\{ x\in P\ |\ a<x<b\right\} $}?
ובכן, לא בדיוק. בואו ניקח קבוצה סדורה פשוטה במיוחד, \L{$\left(\mathbb{N},\le\right)$},
הטבעיים עם יחס הסדר הרגיל. כדי שאוסף כל הקטעים \L{$\left(a,b\right)$}
יהיה בסיס צריך שכל איבר של \L{$\mathbb{N}$} ישתייך לאיבר בסיס כלשהו,
כדי לקיים את התכונה
\begin{itemize}
\item לכל \L{$x\in X$} קיים \L{$B\in\mathcal{B}$} כך ש-\L{$x\in B$}.
\end{itemize}
אבל לאן {\beginL 0\endL} נכנס? הוא הרי לא שייך ל-\L{$\left(0,1\right)$}
כי זו קבוצה שלא מכילה לא את {\beginL 0\endL} ולא את {\beginL 1\endL}
- למעשה, היא ריקה. \L{$\left(0,2\right)$} כוללת את {\beginL 1\endL}
אבל לא את {\beginL 0\endL}, וגם באופן כללי\L{$\left(0,n\right)$}
לא כוללת את {\beginL 0\endL}, וכמובן שאם \L{$0<a$} אז \L{$\left(a,b\right)$}
לא תכלול את {\beginL 0\endL}. אין לנו למשל את \L{$-1$} בתוך הטבעיים
כדי ליצור את הקטע \L{$\left(-1,1\right)$} שיכלול את {\beginL 0\endL}.

אז אנחנו צריכים קצת יותר מזה - אנחנו צריכים איכשהו להתחשב \textbf{בקצוות}
של יחס הסדר. אם ב-\L{$P$} קיים איבר קטן ביותר \L{$a_{0}\in P$} )קטן
ביותר פירושו שלכל \L{$x\in P$} מתקיים \L{$a_{0}\le x$}( אז צריך
להוסיף לבסיס שלנו משהו שמכיל את \L{$a_{0}$}. באופן דומה אם \L{$b_{0}$}
הוא האיבר הגדול ביותר צריך להוסיף משהו שמכיל אותו. אז אולי מספיק להוסיף
את הקבוצות \L{$\left\{ a_{0}\right\} ,\left\{ b_{0}\right\} $}? אפשר
לעשות את זה, אבל נקבל משהו מוזר. בואו נחשוב לרגע על \L{$[0,\infty)$},
כל המספרים הממשיים האי-שליליים. אם אנחנו חושבים על קבוצות פתוחות על
פי המטריקה הרגילה על המרחב הזה, אז כל קבוצה פתוחה שמכילה את {\beginL 0\endL}
היא מהצורה \L{$[0,\varepsilon)$} עבור \L{$\varepsilon>0$} כלשהו;
הקבוצה \L{$\left\{ 0\right\} $} שכוללת רק את {\beginL 0\endL} היא
לא פתוחה בעצמה. זה תואם את האינטואיציה מאחורי קבוצה פתוחה - קבוצה
שמכילה את \textquotedblleft כל הנקודות שנמצאות בסביבה של {\beginL 0\endL}\textquotedblright .
אם רק \L{$\left\{ 0\right\} $} עצמו הוא קבוצה פתוחה, האינטואיציה
היא ש-{\beginL 0\endL} הוא \textquotedblleft מבודד\textquotedblright{}
מיתר המרחב איכשהו. אם כן, אם ב-\L{$P$} קיים איבר קטן ביותר \L{$a_{0}\in P$},
מה ש\textquotedblright טבעי\textquotedblright{} להוסיף לקבוצות הפתוחות
שלנו הוא את כל הקבוצות מהצורה \L{$[a_{0},b)$} עבור כל \L{$a_{0}<b$}
ואם \L{$b_{0}$} הוא איבר גדול ביותר, גם את כל הקבוצות מהצורה \L{$(a,b_{0}]$}
עבור כל \L{$a<b_{0}$}.

אז כדי לסכם, אנחנו מגדירים קבוצה \L{$\mathcal{B}$} שמיועדת להיות
בסיס לטופולוגיה על \L{$P$} וכוללת שלושה סוגי קבוצות )לכל היותר(:
\begin{itemize}
\item כל \L{$\left(a,b\right)$} עבור \L{$a<b$}
\item כל \L{$[a_{0},b)$} עבור \L{$a_{0}<b$} בהינתן שקיים \L{$a_{0}$}
קטן ביותר ב-\L{$P$}
\item כל \L{$(a,b_{0}]$} עבור \L{$a<b_{0}$} בהינתן שקיים \L{$b_{0}$}
גדול ביותר ב-\L{$P$}
\end{itemize}
כמובן, לא מספיק לתת את \L{$\mathcal{B}$} הזו, צריך גם להוכיח שהיא
בסיס. זה שכל איבר של \L{$P$} שייך לאחת מהקבוצות, זה פשוט למדי. ניקח
\L{$x\in P$}. אם הוא לא הגדול ביותר ולא הקטן ביותר קיימים \L{$a,b$}
כך ש-\L{$a<x<b$} ולכן \L{$x\in\left(a,b\right)$}. אם לעומת זאת \L{$x$}
הוא הקטן ביותר אז אם קיים \L{$b$} כך ש-\L{$x<b$}, נקבל ש-\L{$x\in[x,b)$}
ובאופן דומה נטפל גם ב-\L{$x$} קטן ביותר. המקרה החריג היחיד הוא זה
שבו \L{$P$} כוללת רק איבר אחד, אז על הקבוצה הזו לא נגדיר טופולוגיה
בצורה הזו )ממילא יש עליה רק טופולוגיה אחת עם הקבוצות הפתוחות \L{$\emptyset,P$}(.

התכונה השניה יותר מעצבנת כי היא כוללת שלל מקרים לבדיקה. העניין הוא
שחיתוך של שני אברי בסיס יכול לתת או קבוצה ריקה )ואז אין מה להוכיח
עליה( או איבר בסיס חדש )ולכן כל נקודה בו אוטומטית שייכת לאיבר בסיס
שמוכל בחיתוך(, ואני אוותר על להראות את זה באופן מפורש.

אפשר לתאר את טופולוגיית הסדר גם בעזרת תת-בסיס בעזרת מה שנקרא \textbf{קרניים}:
\begin{itemize}
\item \L{$\left(a,\infty\right)=\left\{ x\in P\ |\ a<x\right\} $}
\item \L{$\left(-\infty,b\right)=\left\{ x\in P\ |\ x<b\right\} $}
\end{itemize}
ועכשיו קל לראות ש-\L{$\left(a,\infty\right)\cap\left(-\infty,b\right)=\left(a,b\right)$},
ואם למשל \L{$a_{0}$} הוא האיבר הקטן ביותר אז \L{$\left(-\infty,b\right)=[a_{0},b)$}
, אז אפשר לקבל את כל אברי הבסיס בעזרת הקרניים.

יפה, אז יש לנו עכשיו משפחה שלמה של טופולוגיות - אבל איך הן נראות בפועל,
על קבוצות קונקרטיות? ובכן, כבר ברור לנו איך זה נראה עבור \L{$\mathbb{R}$}
ותת-קבוצות שלו. איזו עוד קבוצה סדורה לינארית אנחנו מכירים? יש אחת
פשוטה במיוחד: \L{$\mathbb{N}=\left\{ 0,1,2,\ldots\right\} $}. למעשה,
מה שהופך את הטבעיים למעניינים הוא שהסדר שלהם הוא \textbf{סדר טוב}
)בכל תת-קבוצה לא ריקה של טבעיים יש איבר קטן ביותר( אבל לא נשתמש בתכונה
הזו כאן.

עבור \L{$\mathbb{N}$}, טופולוגיית הסדר \textbf{כן }נותנת את מה שתיארתי
קודם בתור תופעה מוזרה: ראשית, \L{$[0,1)=\left\{ 0\right\} $}, אז
הסינגלטון שמכיל רק את {\beginL 0\endL} הוא קבוצה פתוחה, אבל גם באופן
כללי - אם \L{$n>0$} אז \L{$\left(n-1,n+1\right)=\left\{ n\right\} $}
ולכן \textbf{כל} סינגלטון הוא קבוצה פתוחה; ומכיוון שאיחוד שרירותי
של קבוצות פתוחות הוא פתוח, זה אומר ש\textbf{כל} קבוצה היא פתוחה. יש
לנו שם לטופולוגיה הזו שבה כל קבוצה היא פתוחה: \textbf{הטופולוגיה הדיסקרטית}.
ה\textquotedblright דיסקרטי\textquotedblright{} מלשון בדיד בשם של הטופולוגיה
מתייחס בדיוק לסיטואציה כמו הנוכחית, שבה \L{$n$} הוא \textquotedblleft בדיד\textquotedblright{}
במובן זה שאין איברים ממש סמוכים אליו.

בואו נעשה את זה מעניין יותר. \L{$\mathbb{N}$} הוא בסך הכל ההתחלה
באוסף המספרים \textbf{הסודרים}, שאני מתאר למשל \href{https://gadial.net/2011/05/30/ordinals_formal_definitions/}{כאן}.
הסודר הבא בתור נקרא \L{$\omega$} )לא חשוב כרגע איך הוא מוגדר פורמלית;
העיקר שיש הגדרה טובה( ואפשר להסתכל על הקבוצה הסדורה \L{$\left\{ 0,1,2,\ldots,\omega\right\} $}.
בקבוצה הזו עם טופולוגיית הסדר עדיין נקבל ש-\L{$\left\{ n\right\} $}
פתוחה לכל \textbf{מספר טבעי}, אבל \L{$\left\{ \omega\right\} $} לא
תהיה פתוחה. מכיוון ש-\L{$\omega$} הוא איבר אחרון, הקבוצות הפתוחות
שהוא שייך אליהן הן \L{$(N,\omega]=\left\{ n+1,n+2,\ldots,\omega\right\} $}
והקבוצות הללו כוללות אינסוף איברים.

מה שמעניין הוא ש-\L{$\omega$} היא \textbf{נקודת הצטברות} של \L{$\mathbb{N}$}
- כל סביבה של \L{$\omega$} כוללת איבר של \L{$\mathbb{N}$}. בנוסף,
\L{$\omega$} היא גם \textbf{נקודת גבול} של \L{$\mathbb{N}$} - היא
הגבול של הסדרה \L{$0,1,2,\ldots$}, כי לכל סביבה של \L{$\omega$}
קיים \L{$N$} כך שהסביבה כוללת את כל \L{$n>N$} )זו בדיוק המשמעות
של \L{$(N,\omega]$}(. ולמה זה מעניין? כי אם ניקח את \L{$\omega$}
ונחליף אותו בסודר אחר שנקרא \L{$\omega_{1}$}, פתאום דברים שעבדו עד
כה יתחילו להישבר. זה מצריך הסבר טיפה יותר זהיר.

\L{$\omega_{1}$} הוא הסודר הראשון שהוא \textbf{לא בן מניה}. פורמלית,
זה אומר שהוא נראה בערך ככה: \L{$\omega_{1}=\left\{ 0,1,2,\ldots,\omega,\omega+1,\ldots\right\} $}
וכן הלאה וכן הלאה - הוא כולל בתוכו \textbf{הרבה} סודרים. אבל כל אחד
מהם הוא בעצמו קבוצה בת מניה.

עכשיו, נסתכל על הקבוצה הסדורה \L{$\omega_{1}+1\triangleq\left\{ 0,1,2,\ldots,\omega_{1}\right\} $}
עם טופולוגיית הסדר. כלומר, לקחנו את \L{$\omega_{1}$} עצמו והוספנו
אותו \textquotedblleft בסוף\textquotedblright{} של עצמו ולתוצאה קראנו
\L{$\omega_{1}+1$} - זה דבר סטנדרטי שעושים עם סודרים. בקבוצה הזו
\L{$\omega_{1}$} הוא איבר אחרון ולכן הקבוצות הפתוחות שמכילות אותו
הן מהצורה \L{$(a,\omega_{1}]$} עבור \L{$a\in\omega_{1}$}. אפשר להוכיח
)לא אעשה את זה כאן( ש-\L{$\omega_{1}$} הוא מה שנקרא \textbf{סודר
גבולי}, כלומר אין סודר \L{$\alpha$} כך ש-\L{$\omega_{1}=\alpha+1$}.
זה מוביל לכך שבכל קבוצה פתוחה מהצורה \L{$(a,\omega_{1}]$} יש אינסוף
איברים - ולכן בדיוק כמו קודם עם \L{$\omega$}, מקבלים ש-\L{$\omega_{1}$}
הוא נקודת הצטברות של הקבוצה \L{$\omega_{1}$} )קודם זה היה פחות מבלבל
כי אמרתי ש-\L{$\omega$} הוא נקודת הצטברות של \L{$\mathbb{N}$}, אבל
הרי ליטרלי \L{$\mathbb{N}=\omega$} אז זה באמת אותו הדבר(.

לאן אני חותר עם זה? לכך ש-\L{$\omega_{1}$} הוא \textbf{נקודת הצטברות}
של \L{$\omega_{1}$} אבל הוא לא \textbf{נקודת גבול} שלה. זו השבירה
הרצינית הראשונה של תוצאה שראינו שמתקיימת במרחבים מטריים: שם הסתמכתי
על תכונה שנקראת \textbf{אקסיומת המניה הראשונה} שמתקיימת במרחבים מטריים
אבל לא באופן כללי.

לא ניסחתי פורמלית את אקסיומת המניה הראשונה עד עכשיו כי הכי קל להסביר
אותה במונחים של \textbf{בסיס}. מה שנקרא \textbf{אקסיומת המניה השניה}
אומר בפשטות שיש למרחב \L{$X$} בסיס \textbf{בן מניה}. אקסיומת המניה
הראשונה היא קצת פחות שאפתנית )וחזקה( וכל מה שהיא דורשת הוא שלמרחב
יהיה \textbf{בסיס מקומי} בן מניה בכל נקודה. מה זה בסיס \textquotedblleft מקומי\textquotedblright{}
בנקודה \L{$x$}? ובכן, בואו נחזור למשהו שאמרתי יותר מוקדם בפוסט: \textquotedblleft היינו
רוצים שבסיס יהיה קבוצה \L{$\mathcal{B}$} כלשהי כך שלכל \L{$A\in\mathcal{T}$}
וכל \L{$a\in A$} יהיה קיים \L{$B\in\mathcal{B}$} כך ש-\L{$a\in B\subseteq A$}\textquotedblright .
אפשר לנסח את התכונה הזו בצורה יותר מילולית - אנחנו רוצים שלכל \L{$a\in X$}
יתקיים שלכל סביבה של \L{$a$}, יש איבר בסיס \L{$B\in\mathcal{B}$}
שמוכל בסביבה הזו, ומהווה בעצמו סביבה של \L{$a$}.

כשעוברים לדבר על \textquotedblleft בסיס מקומי\textquotedblright{} כבר
לא רוצים \L{$\mathcal{B}$} שיעבוד לכל ה-\L{$a$}-ים במרחב בו זמנית
)מה שעלול לאלץ את \L{$\mathcal{B}$} להיות גדול במקרה שבו יש המון
נקודות במרחב( אלא מתחילים עם \L{$a$} ודורשים קיום של קבוצה \L{$\mathcal{B}_{a}$}
של סביבות של \L{$a$} עם התכונה שלכל סביבה \L{$A$} של \L{$a$}, קיים
\L{$B\in\mathcal{B}_{a}$} כך ש-\L{$B\subseteq A$}. \textbf{אקסיומת
המניה הראשונה} אומרת שלכל \L{$a\in X$} קיים בסיס מקומי בן מניה. אין
כאן משהו חדש - בפוסט הקודם זה מה שאני מסתמך עליו, פשוט קצת במובלע.

במקרה של המרחב \L{$\omega_{1}+1$}, לקבוצה \L{$\omega_{1}$} אין בסיס
מקומי בן מניה. אוסף כל אברי הבסיס שמכילים את \L{$\omega_{1}$}הוא
מהצורה \L{$(a,\omega_{1}]$} עבור \L{$a\in\omega_{1}$} ופתחנו בכך
ש-\L{$\omega_{1}$} הוא סודר לא בן מניה, ולכן בפרט האוסף הזה של אברי
הבסיס הוא לא בן מניה. כמובן, זה בפני עצמו לא שולל את האפשרות שאולי
נוכל למצוא בסיס מקומי \textbf{אחר} של \L{$\omega_{1}$} שכן יהיה בן
מניה; בשביל לרסק את התקווה הזו מספיק להראות ש-\L{$\omega_{1}$} הוא
לא נקודת גבול, כפי שהיה אמור לנובע מכך ש-\L{$\omega_{1}$} הוא \textbf{כן}
נקודת הצטברות, אם המרחב שלנו \textbf{כן} היה מקיים את אקסיומת המניה
הראשונה.

ההוכחה תהיה פשוטה, בהינתן שמכירים סודרים; אם לא, אני מודה שזה טיפה
טריקי לעכל את זה עכשיו. מה שנעשה הוא לקחת סדרה כלשהי \L{$\alpha_{1},\alpha_{2},\ldots$}
של איברים של \L{$\omega_{1}$} ולהוכיח שהיא לא מתכנסת אל \L{$\omega_{1}$}.
כל אברי הסדרה הזו הם איברים של \L{$\omega_{1}$} ולכן כולם סודרים
בני מניה. עכשיו נסתכל על הקבוצה

\L{$\left\{ \alpha_{i}\ |\ i\in\mathbb{N}\right\} $}

זו קבוצה של סודרים; אחת מהתכונות הבסיסיות של סודרים היא שאיחוד של
קבוצה של סודרים הוא בעצמו סודר - הסודר הקטן ביותר שגדול או שווה לכל
הסודרים בקבוצה )אני מדבר על זה \href{https://gadial.net/2023/01/14/ordinal_arithmetic/}{כאן}(.
אז מה שיש לנו כאן הוא סודר \L{$\alpha=\bigcup^{\infty}_{i=1}\alpha_{i}$}
שמתקבל בתור \textbf{איחוד בן מניה של קבוצות בנות מניה} ולכן הוא בעצמו
בן מניה, כלומר \L{$\alpha<\omega_{1}$}. נסתכל על הקבוצה הפתוחה \L{$(\alpha,\omega_{1}]$}
- בקבוצה הזו אין \textbf{אף איבר} מהסדרה \L{$\alpha_{1},\alpha_{2},\ldots$}
כי כל האיברים בקבוצה הפתוחה הזו גדולים ממש מ-\L{$\alpha$}. מכאן שהגבול
של הסדרה, אם קיים, לא יכול להיות \L{$\omega_{1}$}, וסיימנו את ההוכחה
ש-\L{$\omega_{1}$} הוא לא נקודת גבול של \L{$\omega_{1}$}.

זה היה קצת... מוזר? ובכן, כן. זה הרעיון - צריך ללכת למקומות מוזרים
כדי להיתקל בטופולוגיות שמתעלות על האינטואיציה שיש לנו ממרחבים מטריים.
זה רק ילך וישתפר בהמשך. לעת עתה בואו נצא להפסקה מבנייה של טופולוגיות
ונעבור לדבר על מושגי הבסיס שעכשיו אפשר לנסח בלשון טופולוגית לחלוטין,
ובראשן - \textbf{פונקציות רציפות}.
\end{document}
